\section{C. The relevant 1-motives and the cuspidal extension}

We keep the arithmetic plus model \(X_n^+\) and its reduced divisor
\(D_n\) of \(3n\) coordinate cusps. All groups and motives in this
part are over characteristic zero. The geometric comparison of B.1
transports the computations on \(x^n+y^n=z^n\) to this model; its
field and its effect on the arithmetic descent are retained.
Write \(J_n=\operatorname{Pic}^0(X_n^+)\) and
\(\mathcal J_n=\operatorname{Pic}^0(X_n^+,D_n)\), the degree-zero
generalized Jacobian with reduced modulus \(D_n\).

\begin{statement}{C.1}{Known}{The actual Fermat 1-motive and its realizations}
Put
\[
 L_n=\operatorname{Div}^0_{D_n}(X_{n,\overline{\mathbb Q}}^+),\qquad
 \mathcal P_n=[L_n\xrightarrow{D\mapsto\mathcal O(D)}J_n].
\]
This is a Deligne 1-motive, with its natural descent over \(\mathbb Q\).
Its Cartier dual, using the principal polarization of \(J_n\), is
\[
 \mathcal A_n=[0\longrightarrow\mathcal J_n].
\]
There is an exact sequence of semiabelian varieties
\[
 1\longrightarrow
 T_{D_n}:=(\operatorname{Res}_{D_n/\mathbb Q}\mathbb G_m)/\mathbb G_m
 \longrightarrow\mathcal J_n\longrightarrow J_n\longrightarrow1.
                                                               \tag{C.1}
\]
The character lattice of \(T_{D_n}\) is \(L_n\), and its cocharacter
lattice is \(\mathbb Z[D_n]/\mathbb Z\sum_{P\in D_n}P\).
The Betti and Tate realizations of \(\mathcal A_n\) are
\[
 T_{\mathbb Z}(\mathcal A_n)=H_1(U_n^+(\mathbb C),\mathbb Z),\qquad
 T_\ell(\mathcal A_n)=
 H_1^{\mathrm{et}}(U_{n,\overline{\mathbb Q}}^+,\mathbb Z_\ell).
                                                               \tag{C.2}
\]
Its de Rham realization is the dual of
\(H^1_{\mathrm{dR}}(U_n^+/\mathbb Q)\).
Its weights are \(-2\), with rank \(3n-1\), and \(-1\), with
rank \((n-1)(n-2)\); its total rank is \(n^2+1\).
\emph{Source: Deligne, \cite[10.1.3--10.1.8, 10.2.10--10.2.11,
10.3.4--10.3.8]{DeligneHodgeIII}.}
\end{statement}
\begin{proof}
We spell out the reduction to the cited construction. A point of the
generalized Picard group is a line bundle on \(X_n^+\), together with
a trivialization along \(D_n\), modulo scalar automorphisms of the
line bundle. The fibre over a fixed degree-zero line bundle is thus
the product of the nonzero fibres at the cusps modulo simultaneous
scaling. This proves (C.1) and both of its lattice descriptions.

Deligne's definition 10.3.4 assigns to a smooth open curve
\(X-D\) the complex
\([\ker(\mathbb Z[D]\to\mathbb Z)\to\operatorname{Pic}^0(X)]\).
For \(X=X_n^+\), this is precisely \(\mathcal P_n\).
Construction 10.3.6 identifies its \(m\)-torsion realization with
\(H^1(U_n^+,\mu_m)\): a representative is a pair
\((\mathcal L,D)\) with \(\mathcal L^{\otimes m}\simeq\mathcal O(D)\),
modulo pairs \((\mathcal O(E),mE)\). This is the same complex
defining the finite realization of \([L_n\to J_n]\).
Thus the comparison includes the cusp lattice and has no missing
toric summand. Construction 10.2.11 dualizes this complex to the
extension by \(\operatorname{Hom}(L_n,\mathbb G_m)\) whose class
evaluated at \(D\) is \(\mathcal O(D)\); the line-bundle description
above identifies this extension with \(\mathcal J_n\).
Dualizing the finite realizations with values in \(\mu_m\)
gives \(H_1(U_n^+,\mathbb Z/m)\), and their inverse limits give (C.2).

The Betti construction 10.1.3 and the comparison 10.3.8 give the
integral homology realization; 10.1.7--10.1.8 and 10.3.9 give the
de Rham comparison. Concretely, its dual Hodge filtration has
\[
 F^1H^1_{\mathrm{dR}}(U_n^+)
   =H^0(X_n^+,\Omega^1_{X_n^+}(\log D_n)),
 \quad \dim F^1=g_n+3n-1.
\]
The torus in (C.1) contributes weight \(-2\) and the abelian
variety weight \(-1\). Their ranks are \(3n-1\) and \(2g_n\),
respectively. B.6 gives \(2g_n=(n-1)(n-2)\).
Every construction used here is functorial for the deck and
permutation automorphisms, which proves compatibility of all
the realizations with those actions.
\end{proof}

\begin{statement}{C.2}{Refuted}{The rank-two identification in the lead}
The assertion that the rank-two module \(\widehat M(1)\) of \S51,
or its reduction modulo a nontrivial integer \(m\), is the full
Tate realization of the generalized Jacobian of \(X_n^+\) with its
cusps is false for \(n>1\).
Its exact replacement is
\[
 \widehat M(1)=T_\bullet(T),\qquad
 T=(\mathbb G_m^3)/\mathbb G_m=\mathcal J_1,
                                                               \tag{C.3}
\]
and the covering pushforward has the exact realization
\[
 0\longrightarrow
 (R_n/\mathbb Z\mathcal N_n)\otimes\widehat{\mathbb Z}
 \longrightarrow T_\bullet(\mathcal J_n)
 \xrightarrow{q_{n*}}n\widehat M(1)\longrightarrow0.             \tag{C.4}
\]
The Galois structure on the first term is the one transported by its
inclusion into the geometric homology in (C.2).
At level \(n=2\), the ranks in (C.4) are \(3,5,2\).
\emph{Manuscript: \S51.11, folio \(697\), scan 99, (228);
\S51.13--14, folios 719--720, scans 143--146, (299)--(305).}
\end{statement}
\begin{proof}
At \(n=1\) the compact curve is a projective line, so its ordinary
Jacobian vanishes. Formula (C.1) then gives (C.3), with cocharacter
lattice \(M\). Abel--Jacobi sends a peripheral loop to the corresponding
cocharacter; hence its Tate realization is the geometric abelianization
used in B.2 and B.14.
For arbitrary \(n\), functoriality of the generalized Jacobian gives
the norm map induced by \(q_n\). On the Betti realization it is
the map \(K_n^{\rm ab}\to F^{\rm ab}\). B.5 computes its kernel
and its image exactly. Tensoring (B.2) with the flat ring
\(\widehat{\mathbb Z}\) and using (C.2) proves (C.4).
This is an explicit comparison morphism; at level two its matrix
and its Smith entries \(2,2\) are supplied in B.7.

The geometric module in the middle has rank \(n^2+1\); the proposed
module has rank two. Modulo \(m>1\), their orders are
\(m^{n^2+1}\) and \(m^2\), respectively. Thus even an abstract
isomorphism of the proposed finite realizations fails, with
\(n=m=2\) giving the finite orders \(32\) and \(4\).
The finite order comparison is checked in
\texttt{code/b\_freeness.py}. The norm comparison (C.4) gives the actual relationship that
replaces that proposed identification.
\end{proof}

\begin{statement}{C.3}{Proved}{Complex multiplication with its actual ring map}
For \(n\ge3\), over a field containing \(\mu_n\), the deck action gives a homomorphism
\[
 R_n\longrightarrow\operatorname{End}(J_n)
\]
whose kernel is \(I_n\) of B.6. The induced commutative semisimple
algebra
\[
 \mathcal E_n=\mathbb Q[A_n]/(I_n\otimes\mathbb Q)
 \ \subset\ \operatorname{End}^0(J_n)
\]
has dimension \((n-1)(n-2)=2\dim J_n\). Its factors are cyclotomic
fields indexed by the Galois orbits of characters
\((a,b)\) with \(a,b,a+b\ne0\pmod n\). This is the precise
complex-multiplication algebra supplied by the Fermat automorphisms.
The action on \(\mathcal J_n\) restricts to the permutation action
on the torus lattice in (C.1).
\end{statement}
\begin{proof}
Automorphisms of the curve act on divisor classes and therefore on
the Jacobian. An endomorphism of a complex abelian variety is zero
as soon as it is zero on integral first homology: after lifting to
the complex universal cover it is a complex linear map, and the
integral lattice spans that cover over \(\mathbb R\).
B.6 identifies this homology with \(R_n/I_n\), so the kernel is
exactly \(I_n\). The rational group algebra of a finite abelian
group is a product of number fields, and its character
decomposition selects exactly the characters in B.6. Each has
multiplicity one. The dimension assertion follows by counting them.
Complex conjugation pairs \((a,b)\) with \((-a,-b)\);
none of the selected characters is real: a real character would
have \(a,b\) of order at most two and would force one of
\(a,b,a+b\) to vanish. Thus the selected cyclotomic factors are
CM fields. This proves the full degree required for complex
multiplication.
The line-bundle-with-trivialization description of \(\mathcal J_n\)
shows directly that deck and coordinate automorphisms permute the
cuspidal fibres, giving the torus action claimed. The notation
\(\mathbb Z[\mu_n]^2\) in Lead 3 is therefore replaced by this
explicit group-ring homomorphism and its kernel, with no
interpretation of that ambiguous expression left implicit.
\end{proof}

\begin{statement}{C.4}{Proved}{Exact order of cuspidal differences}
For every \(n\ge3\) and every two distinct geometric coordinate
cusps \(P,Q\) of \(X_n^+\), the class \([P-Q]\in J_n\) has exact
order \(n\). Consequently the subgroup generated by all
cuspidal degree-zero divisor classes has exponent exactly \(n\).
At \(n=1,2\) that subgroup is zero.
\end{statement}
\begin{proof}
Let \(\ell_Q=0\) be the tangent line at a coordinate cusp \(Q\).
For example, at \(Q=[0:q:1]\), with \(q^n=-1\), its equation is
\(y-qz=0\). Substitution into \(x^n+y^n+z^n=0\) gives \(x^n=0\);
the intersection divisor is exactly \(nQ\). The same holds at
every coordinate cusp. Therefore
\[
 \operatorname{div}(\ell_P/\ell_Q)=n(P-Q),                    \tag{C.5}
\]
and the order divides \(n\).

We prove exactness, including the even levels. At \(Q=[0:q:1]\),
the functions \(z/\ell_Q\) and \(x/\ell_Q\) have their only poles
at \(Q\), of orders \(n\) and \(n-1\), respectively. Thus the pole
semigroup \(H(Q)\) contains \(\langle n-1,n\rangle\).
The latter numerical semigroup has exactly
\((n-1)(n-2)/2\) positive gaps: the least element in residue
class \(r\pmod{n-1}\) is \(rn\), for \(0\le r<n-1\), and counting
the smaller positive members of those residue classes gives this
number. Riemann--Roch gives exactly \(g_n\) gaps for \(H(Q)\):
for \(a\ge2g_n\), \(\ell(aQ)=a+1-g_n\), and every successive
dimension increase is zero or one. B.6 gives
\(g_n=(n-1)(n-2)/2\). Since the contained semigroup already has
all \(g_n\) gaps, it equals \(H(Q)\).

A proper positive divisor \(d\) of \(n\) satisfies
\(d\le n/2<n-1\), hence is a gap. A principal divisor
\(d(P-Q)\) would supply a nonconstant function with its only
pole at \(Q\), of order \(d\), a contradiction. The exact order
is therefore \(n\). The genus-zero cases have zero Jacobian.
This gives a proof for every \(n\), without a parity restriction.
\end{proof}

\begin{statement}{C.5}{Proved}{The extension class and its exact isogeny exponent}
The class of (C.1) is the homomorphism
\[
 L_n\longrightarrow J_n^\vee,\qquad D\longmapsto[D],
                                                               \tag{C.6}
\]
under the principal polarization and the extension--character
correspondence
\(\operatorname{Ext}^1(J_n,T_{D_n})
 \simeq\operatorname{Hom}(L_n,J_n^\vee)\).
For every \(n\ge3\), this class has exact order \(n\).
Its pullback by \([n]:J_n\to J_n\) is split, and the resulting
isogeny
\[
 T_{D_n}\times J_n\longrightarrow\mathcal J_n
\]
has kernel isomorphic to \(J_n[n]\) and degree \(n^{2g_n}\).
The splitting of that pullback descends to \(\mathbb Q\).
There is no single positive integer that kills these extension
classes at every Fermat level.
\end{statement}
\begin{proof}
Push (C.1) out by the torus character associated with
\(D=\sum a_PP\). The fibre of the resulting \(\mathbb G_m\)-extension
over a line bundle \(\mathcal L\) is the nonzero part of
\(\bigotimes_P\mathcal L(P)^{\otimes a_P}\).
This is the restriction of the Poincare bundle corresponding
to \([D]\in J_n^\vee\); it proves (C.6).
Conversely, these character pushouts determine a torus extension,
because a split torus is a product of copies of \(\mathbb G_m\).
Descent retains exactly the Galois-equivariant homomorphisms in
(C.6); morphisms from an abelian variety to a torus vanish,
so there is no ambiguity in descending a splitting.
This is also the Cartier-duality description in
\cite[10.2.11--10.2.12]{DeligneHodgeIII}.

C.4 says exactly that (C.6) is killed by \(n\), and that no
proper divisor of \(n\) kills it. Pullback by the multiplication
map \([m]\) multiplies the associated point of \(J_n^\vee\)
by \(m\), as is immediate from the Poincare bundle. Thus
\([n]^*(\text{C.1})\) is the trivial torus extension.
Projection from that pullback to \(\mathcal J_n\) has kernel
\(J_n[n]\), proving the degree and the isogeny statement.
A splitting over \(\overline{\mathbb Q}\) is unique because
\(\operatorname{Hom}(J_n,T_{D_n})=0\), so it is Galois-invariant
and descends. Finally C.4 gives exponents \(n\) for unbounded
\(n\), proving the last assertion.

This statement describes the order in the algebraic extension
group. Its realizations give a splitting of the weight
filtration in the isogeny category, and an integral comparison
whose finite defect is controlled by \(n\). It does not replace
that specified extension group by an unspecified Galois
cohomology group.
\end{proof}

\begin{statement}{C.6}{Known}{Rohrlich's theorem and the requested computations}
For \(n\ge3\), the geometric cuspidal subgroup is
\[
 C_n\simeq
 \begin{cases}
  (\mathbb Z/n)^{3n-7},&n\text{ odd},\\
  (\mathbb Z/n)^{3n-7}\times\mathbb Z/2,&n\text{ even}.
 \end{cases}
\]
This is \cite[Theorem 1, p.~101]{Rohrlich1977},
\emph{Points at infinity on the Fermat curves}. We opened the original
theorem and its definitions on p.~95. The odd case is also reproduced
with an independent proof in \cite[Theorem 1, p.~28-02]{Velu1979}.
For the requested levels the exact orders of the extension class
are
\[
\begin{array}{c|rrrr}
n&3&4&5&7\\ \hline
g_n&1&3&6&15\\
\operatorname{rk}T_{D_n}&8&11&14&20\\
\operatorname{rk}T_\bullet(\mathcal J_n)&10&17&26&50\\
\operatorname{ord}[(\mathrm{C.1})]&3&4&5&7
\end{array}
\]
At \(n=3,5,7\), the computed nontrivial Smith factors of the
checked cuspidal presentation are respectively
\(3^2,\ 5^8,\ 7^{14}\).
\end{statement}
\begin{proof}
Rohrlich defines \(\mathcal D^\infty\) to be the degree-zero divisors
supported on the \(3n\) coordinate cusps of \(x^n+y^n=z^n\), and
\(\mathcal F^\infty\) to be the principal divisors in this lattice.
The geometric isomorphism \(\iota_n\) of B.1 identifies their quotient
with the image of \(L_n\to J_n\). By (C.6), that is precisely the image
of the homomorphism classifying the generalized-Jacobian extension.
The exponent of the displayed group is \(n\), also in the even case
because \(2\mid n\). Therefore multiplication by \(n\) kills the
extension class, and no smaller positive integer does. C.5 then
constructs the splitting isogeny, over \(\mathbb Q\) itself.
This is the reduction from the original theorem to the asserted
weight-extension statement.

For independent integral computations at the requested odd levels,
choose a cusp
\(P_*\), and identify the degree-zero cusp lattice with the
free lattice on \(P-P_*\), for \(P\ne P_*\).
Velu's proof determines its principal-divisor sublattice modulo
\(n\), with six relation classes. In the standard ordering
\(a_i,b_i,c_i\), these can be written
\[
\begin{gathered}
 \sum_i(a_i-P_*),\quad \sum_i(b_i-P_*),\quad
 \sum_i(c_i-P_*),\\
 \sum_i i(a_i-b_i),\quad \sum_i i(a_i-c_i),\quad
 \sum_i i^2(a_i+b_i+c_i-3P_*).
\end{gathered}
\]
The same displayed presentation appears in
\cite[Theorem 29]{BanerjeeMerel2022}, which was also opened.
Together with \(n\) times each basis vector it is the relation
matrix used by \texttt{code/c\_cuspidal\_exponent.py}.
Its exact integer Smith form gives the three listed odd-level
groups, hence their exponents. Formula (C.6) identifies the image
of the extension-class homomorphism with that same cuspidal group,
so the exponent is again the order of the extension class.

For \(n=4\) the program does not apply an odd-level presentation.
It checks instead the independent exact certificate of C.4:
the tangent divisor is \(4(P-Q)\), the pole semigroup
\(\langle3,4\rangle\) has gaps \(1,2,5\), and the proper divisors
\(1,2\) of \(4\) are gaps. Therefore the class has exact order
four. The same certificate is checked for all four requested
levels and, as an additional arithmetic verification, for every
\(3\le n\le32\). The proof of C.4 supplies its validity for all
levels.
\end{proof}

\begin{statement}{C.7}{Proved}{The toric motive actually constructing the tower}
Let \(T=\mathbb G_m^3/\mathbb G_m\), and let
\(\psi:U_1^+\hookrightarrow T\) be the map of B.1. There is a
Cartesian square over \(\mathbb Q\)
\[
\begin{array}{ccc}
 U_n^+&\longrightarrow&T\\
 {\scriptstyle q_n}\downarrow&&\downarrow{\scriptstyle[n]}\\
 U_1^+&\xrightarrow{\ \psi\ }&T.
\end{array}                                                   \tag{C.7}
\]
For a point \(e\in U_1^+(k)\), the 1-motive
\[
 \mathcal K_e=[\mathbb Z\xrightarrow{1\mapsto\psi(e)}T]
\]
has finite realizations
\[
 0\to T[n]\to T_{\mathbb Z/n}(\mathcal K_e)\to\mathbb Z/n\to0,
                                                               \tag{C.8}
\]
whose fibre over \(1\) is \(q_n^{-1}(e)\).
The compatible fibres give exactly the Fermat torsor and its
affine Galois action of B.14. Thus the rank-two module of \S51
is the weight-\(-2\) part of this toric Kummer construction.
\emph{Manuscript: \S51.13--14, folios 719--724, scans 143--153,
(299)--(323), especially (317) and (323).}
\end{statement}
\begin{proof}
A point \([a_0:a_1:a_\infty]\) of \(T\) lies over
\([-e:e-1:1]\) under \([n]\) precisely when its \(n\)-th powers
have that projective ratio. Their sum is zero, which is the
equation of \(U_n^+\). This proves the Cartesian square
scheme-theoretically by its equations on a torus chart.
Deligne's construction 10.1.5 represents
\(T_{\mathbb Z/n}([\mathbb Z\to T])\) by pairs
\((a,t)\), with \(t^n=\psi(e)^a\), modulo
\((a,t)\sim(a+nb,t\psi(e)^b)\). Projection to \(a\bmod n\)
gives (C.8), and the fibre at \(a=1\) is the torsor of
\(n\)-th roots of \(\psi(e)\). By (C.7) it is the stated
Fermat fibre. All formulas commute with automorphisms of \(k\)
and with coordinate permutations, proving the claimed action
and its compatibility with the representation in B.14.

This construction has weights \(0,-2\), with quotient
\(\mathbb Z\) and toric Tate kernel of rank two.
The generalized-Jacobian construction in C.1 is connected to it
by the norm morphism (C.4); that map records the full additional
kernel at level \(n\).
\end{proof}

\begin{statement}{C.8}{Refuted}{Lead 3 and the reading of the manuscript}
Lead 3's proposed identification of \S51.13--14 with a splitting
of the generalized Jacobian of every Fermat curve is false.
The source's explicit construction is (C.7) and the Kummer
fibre (C.8). Its further elliptic-curve construction starts
with an elliptic curve \(E/S\) and the three nonzero
two-torsion points \(D=E[2]\setminus\{0\}\), and has exact sequence
\[
 1\to \mathbb G_m\otimes
       (p_*\mathbb Z_D/\mathbb Z_S)
   \to\mathcal J_{E,D}\to E\to1.
                                                               \tag{C.9}
\]
The Tate kernel of (C.9) has rank two, as written in (303).
It is identified with the torus of the three-point quotient
\(E/\{\pm1\}\) punctured at the images of \(D\).
The weight-extension result for the Fermat Jacobians is C.5,
with exact and unbounded exponents \(n\); the rank-two group
extension of B.12 has zero class and unique splitting conjugacy
class.
\emph{Manuscript: \S51.13, folios 719--720, scans 143--145,
(299)--(304); \S51.14, folios \(720'\)--724, scans 146--153,
(305)--(323).}
\end{statement}
\begin{proof}
The divisor used in scan 144 is explicitly the degree-three
subscheme of nonzero two-torsion points on the elliptic curve,
and (302) has quotient that elliptic curve. Applying the
line-bundle construction of C.1 to \((E,D)\) proves (C.9).
The three-point quotient map \(E\to E/\{\pm1\}\) identifies the
two degree-three etale divisors. Hence it identifies their
cocharacter lattices \(p_*\mathbb Z_D/\mathbb Z_S\) and their
associated tori. On the generalized Jacobians this torus isomorphism
is the pullback of a trivial line bundle with prescribed nonzero
values at the three marked points; the three values are carried
to the corresponding points of \(D\), modulo common scaling.
This is the explicit comparison explaining
the repeated rank-two Tate module. The degree-three divisor
is etale over bases on which \(2\) is invertible, precisely
the qualification made at scan 152.

The other torus comparison written at scan 147 is also explicit:
\[
 [x_0:x_1:x_\infty]\longmapsto
 (x_1/x_\infty,x_\infty/x_0,x_0/x_1)
 \in\ker(\mathbb G_m^3\xrightarrow{\mathrm{product}}\mathbb G_m).
\]
It is an isomorphism: its first two coordinates recover
\(x_1/x_\infty\) and \(x_0/x_\infty\).
For an even permutation it commutes with coordinate permutation;
for an odd permutation its value is the inverse of the permuted
value. This proves the manuscript's marginal qualification that
it is not equivariant for the two untwisted full \(S_3\)-actions.
The \(S_3\)-module used in B.10 is always the quotient lattice,
so no product-one torus action is substituted in the vanishing proof.

At scans 151 and 153 the manuscript supplies \(\psi\) and the
fibre of the universal multiplication tower at the marked
section \(e\). They are exactly the equations proved in C.7.
For the full Fermat generalized Jacobian, the torus has rank
\(3n-1\) and the comparison with this rank-two torus is the
noninjective norm morphism (C.4). Thus the proposed equality
fails by that explicit morphism and its nonzero kernel already
at \(n=2\). The exponents in C.5 further show that a fixed
bounded torsion correction cannot describe all levels.

The supplied table of contents retains the unread initial
letters in \emph{[??]tivisation}; the body heading at folio 708, scan 121,
reads \emph{Retour sur la tour de Fermat}. The mathematics
settles the interpretation just described. It does not settle
the missing letters of the table-of-contents word.
\end{proof}
